\documentclass[exercice,a4]{thedocument}%
\usepackage{texmaths-tikz}%
\usepackage{pstricks}%
\usepackage{pst-bezier}%
\usepackage{pstricks-add}%
\usepackage{pst-eucl}%
\usepackage{pst-func}%
\usepackage{pst-tree}%
\usepackage{pst-circ}%
\usepackage{tikz}%
\usepackage{pgfplots}%
\usepackage{tkz-tab}%
\usepackage{tikzpagenodes}%
\usepackage{tabularray}%
\usepackage{tabularx}%
\usepackage{multirow}%
\usepackage{multicol}%
\usepackage{scratch3}%
\usepackage{enumitem}%
\usepackage[np]{numprint}%
\usepackage{pas-tableur}%
\usepackage{listings}%
\usepackage{graphicx}%
\usepackage{adjustbox}%
\usepackage{needspace}%
\usepackage{xcolor}%
\usepackage{eurosym}%
\usepackage{diagbox}%
\usepackage{colortbl}%
\usepackage{multido}%
\usepackage{icomma}%
\usepackage[french]{babel}%
\startdatakeys
    \source{26GENMATMEAG3}
    \BO{2026}
    \cycle{4}
    \competencesDuSocle{CA,RA,CO}
    \objectifsApprentissage{C4A1aa,C4A2ad,C4B2cb,C4D2aa,C4D3ac}
    \motsCles{Pythagore,périmètre,aire}
\enddatakeys
\definecolor{gris}{gray}{0.5}%
\providecommand{\wideparen}[1]{\overset{\frown}{#1}}%
\begin{document}%
\providecommand{\degres}{^\circ}%
Le schéma ci-dessous représente un terrain SMVU. Il n'est pas à l'échelle.

%Début Tikz
\begin{center}
\begin{tikzpicture}[x=0.26cm,y=0.26cm,line cap=round, line join=round]
	\coordinate (S) at (0,12);
	\coordinate (T) at (27,12);
	\coordinate (M) at (36,12);
	\coordinate (U) at (0,0);
	\coordinate (V) at (27,0);
	% angles droits
	\draw[line width=1.2pt,red,fill=red!30]
		(S) rectangle ++(1.5,-1.5)
		(T) rectangle ++(1.5,-1.5)
		(U) rectangle ++(1.5,1.5)
		(V) rectangle ++(-1.5,1.5);
	\draw[line width=1.2pt] (U) -- (S) -- (M) -- (V) -- cycle;
	\draw[line width=1pt,dashed] (T) -- (V);
	% ouverture du portail
	\draw[line width=1pt] (10,-0.6) -- (10,0.6) (13,-0.6) -- (13,0.6);
	\draw[line width=1.5pt,white] (10.1,0) -- (12.9,0);
	\draw[line width=1pt,densely dotted] (10,0) -- (13,0);
	% cotes
	\draw[<->, line width=0.8pt, > = stealth] (0,15) -- (36,15) node[pos=0.5,above]{36 m};
	\draw[<->, line width=0.8pt, > = stealth] (-3.5,0) -- (-3.5,12) node[pos=0.5,left]{12 m};
	\draw[<->, line width=0.8pt, > = stealth] (0,-3) -- (27,-3) node[pos=0.5,below]{27 m};
	\draw[line width=0.4pt,gray,dashed] (0,12) -- (0,15.5) (36,12) -- (36,15.5)
		(0,12) -- (-4,12) (0,0) -- (-4,0)
		(0,0) -- (0,-3.5) (27,0) -- (27,-3.5);
	\node[above left] at (S) {S};
	\node[above left] at (T) {T};
	\node[right] at (M) {M};
	\node[below left] at (U) {U};
	\node[below right] at (V) {V};
\end{tikzpicture}
\end{center}
%Fin Tikz

On souhaite poser du gazon sur la partie rectangulaire STVU et cultiver un potager sur la partie triangulaire TMV.

On désire également clôturer ce terrain SMVU avec du grillage en laissant une ouverture de 3~m pour le futur portail qui sera placé sur le côté [UV].

\begin{enumerate}
	\item Justifier que la longueur TM est égale à 9~m.
	\item Montrer que la longueur MV est égale à 15~m.
	\item Quelle longueur de grillage doit-on prévoir ?
	\item Montrer que la surface de la partie du terrain STVU recouverte de gazon est égale à 324~m$^2$.
	\item Le cadre ci-dessous donne les informations que l'on peut lire sur l'étiquette d'un rouleau de gazon.

	\begin{center}
	\begin{tblr}{colspec={|l|},row{1}={font=\bfseries,halign=c}}
		\hline
		Caractéristiques d'un rouleau une fois déroulé \\
		Forme : rectangulaire \\
		Longueur : 4~m \\
		Largeur : 2~m \\
		Épaisseur : 30~mm \\
		Prix : 86~\euro \\ \hline
	\end{tblr}
	\end{center}

	On veut recouvrir de gazon le rectangle STVU.

	Un budget de \np{3500}~\euro{} est-il suffisant ?\\
	\textbf{Détailler la réponse en précisant la démarche.}
\end{enumerate}

\needspace{10\baselineskip}
\medskip
\showthesource
\end{document}
